\chapter{Discussion, Conclusions and Future Scope}\label{ch:concl}

\section{Discussion and Current Findings}

\paragraph{Accuracy per LUT.} The 50-LUT DWN gives the best accuracy per unit of hardware among the models compared:
an event-F1 of 0.211 with 2.4 false alarms per day at about 960 estimated LUTs and no multipliers, against 0.266 for
TreeLUT at about 6.4k LUTs and 0.292 for LightGBM, which is not a hardware design. The INT8 MLP, the conventional way
to put a neural network on a chip, is worse on every metric while needing 9,600 multiplications per decision. Within
the DWN family, size did not matter between 50 and 200 LUTs, while the choice of input features and the encoder
resolution did.

\paragraph{The encoder dominates the cost.} About 850 of the 960 estimated LUTs are comparators of the thermometer
encoder; the network itself is 50 LUTs. This agrees with recent measurements of DWN hardware \cite{mecik2025thermo} and
makes the encoder (extension N2) the main lever for reducing the hardware further.

\paragraph{False alarms look like seizures to the network.} The interpretability study showed that the network reads
almost the same inputs during true and false alarms. It also exposed a weakness of the first feature selection, which
ignored the heart signal; fixing it was the largest single improvement of Phase I. Inspecting the trained LUTs was
therefore more useful than tuning the network blindly.

\paragraph{Tonic-clonic against focal seizures.} The detector finds 88\,\% of tonic-clonic seizures but 8--17\,\% of
focal seizures, and loses most of its performance without ECG. This matches the literature, where wearable detection is
established only for tonic-clonic seizures \cite{beniczky2021guideline}, and it shows where the remaining work lies:
focal seizures and robustness to missing sensors.

\paragraph{Reproducibility.} The published ChronoNet baseline could only be partly reproduced, and its released code
selects the model on the test patients. Keeping our own test patients untouched until the design is frozen, and
reporting a clean version of the baseline, is therefore essential for an honest final comparison.

\section{Current Progress and Future Scope}

Table~\ref{tab:phase1-months} summarises the work carried out in the two months of Phase I, and
Table~\ref{tab:status} the status of every activity at the end of Phase I.

\begin{table}[htbp]
\centering\small
\caption{Work carried out during Phase I.}\label{tab:phase1-months}
\begin{tabularx}{\textwidth}{@{}P{30mm}L@{}}
\toprule
\textbf{Period} & \textbf{Activities}\\
\midrule
Month 1 (September 2026) & Selection of the topic and definition of the problem with the guide; literature study on
wearable seizure detection and LUT-based neural networks; study of the SeizeIT2 dataset, its baselines and its
evaluation; setting up the software tools and data access; first experiments, whose results were weaker than those
reported here.\\
Month 2 (October -- November 2026) & Complete audit of the dataset; reproduction of the ChronoNet baseline (official
and clean); integer feature pipeline and fixed-point specification; cross-validated comparison of nine model families;
design study of the DWN and the improved feature selection; first extensions (N1, N3, N7); hardware handoff tools;
this report.\\
\bottomrule
\end{tabularx}
\end{table}

\begin{xltabular}{\textwidth}{@{}P{55mm}P{26mm}L@{}}
\caption{Status of the project activities at the end of Phase I.}\label{tab:status}\\
\toprule
\textbf{Activity} & \textbf{Status} & \textbf{Remarks}\\
\midrule
\endfirsthead
\toprule
\textbf{Activity} & \textbf{Status} & \textbf{Remarks}\\
\midrule
\endhead
\bottomrule
\endlastfoot
\small Literature review & \small Completed & \small 79 references in five areas; research gaps identified\\
\small Dataset access and audit & \small Completed & \small All 11,009 files measured; differences from the
documentation recorded\\
\small Reproduction of ChronoNet & \small Completed & \small Partly reproduced; the clean run is the reference for the
final comparison\\
\small Reproduction of the SVM baseline & \small Pending & \small Needs the 42-feature definition of the original
study\\
\small Integer feature pipeline and fixed-point specification & \small Completed & \small 117 features; no measurable
accuracy loss\\
\small Comparison of model families & \small Completed for nine families & \small LogicNets-style network, GRU and
raw-EEG CNN/TCN running\\
\small DWN design study & \small Completed up to 200 LUTs & \small 500--2,000 LUTs and further ablations running on the
GPU\\
\small Extensions N1, N3, N7 & \small Completed & \small None adopted alone; N7 led to the improved feature selection\\
\small Extensions N2, N4, N5, N6 & \small In progress & \small 16 variants running; missing-sensor tests to follow\\
\small Hardware handoff tools (export, test vectors, reference Verilog) & \small Completed and tested & \small Run on
the final model in Phase II\\
\small Final evaluation on the 29 test patients & \small Planned & \small After every design choice is frozen; five
seeds per model\\
\small FPGA implementation & \small Planned & \small Synthesis, measurements and demonstration in Phase II\\
\end{xltabular}

The remaining months will concentrate on the open questions of Phase I. The extensions will be completed and the
winners combined; the large LUT networks and the raw-EEG references will close the model comparison; the design will
then be frozen and evaluated once on the 29 test patients. The final DWN will be exported as truth tables with test
vectors, implemented on a low-power FPGA and measured against the INT8 MLP and TreeLUT on the same device.
Table~\ref{tab:plan} gives the month-wise plan.

\begin{table}[htbp]
\centering\small
\caption{Month-wise plan for the remaining project period.}\label{tab:plan}
\begin{tabularx}{\textwidth}{@{}P{30mm}LP{42mm}@{}}
\toprule
\textbf{Month} & \textbf{Work} & \textbf{Expected deliverable}\\
\midrule
November -- December 2026 & Remaining extensions (N2, N4, N5, N6) and their combination; large DWNs, GRU and raw-EEG
references; design frozen & Ablation table; final DWN design\\
December 2026 & Final evaluation on the 29 test patients with five seeds; hardware handoff package & Test-set results;
handoff package with verified reference Verilog\\
January 2027 & FPGA synthesis and place-and-route; conference paper draft (EMBC 2027, full papers reportedly due 24
January 2027) & Resource and timing report; paper draft\\
February 2027 & Measured LUTs, power and latency; INT8 MLP and TreeLUT on the same device & Hardware comparison table\\
March 2027 & On-board demonstration with recordings streamed through the design; documentation & Demonstration\\
April 2027 & Final report and presentation & Final report\\
\bottomrule
\end{tabularx}
\end{table}

\section{Conclusions}

Phase I established the foundation of the project and its first results:
\begin{enumerate}
\item The SeizeIT2 dataset was audited completely; its sampling rate, seizure count and labels differ from the
  documentation, and the design follows the files.
\item The published ChronoNet baseline was reproduced only partly, and a clean version without test information was
  trained as the reference for the final comparison.
\item A bit-exact integer feature pipeline for behind-the-ear EEG, ECG, EMG and motion loses no measurable accuracy
  against floating point.
\item A 50-LUT Differentiable Weightless Neural Network reaches an event-F1 of 0.211 with 2.4 false alarms per day on
  cross-validation, at about 960 estimated LUTs and no multipliers: seven to eight times smaller than the other LUT
  models and better than an INT8 network, though still below gradient-boosted trees.
\item Interpretability of the LUT network led to the largest improvement of the phase, and the thermometer encoder was
  identified as the dominant hardware cost.
\item No LUT-native seizure detector exists in the literature, and no LUT network fuses modalities or handles missing
  sensors; the project addresses both.
\end{enumerate}
Phase II will complete the extensions, evaluate the frozen design once on the 29 held-out patients, and implement and
measure the detector on an FPGA.
